\scanpage{65}
\mspage{- 8 -}

\margin{\handwritten{$\checkmark$}}\[ \cup C_n : H^q(X, Z_n) \xrightarrow{\unsure{\sim}} H^{q+2d}_X (E, \mu_n^{\otimes d}) . \]

When $X^\prime = G_k$ , the $\mu_n$ in this formula refers to the abelian
group $\mu_{n,k}$ of nth roots of unity in $k$ . But the map $k^* \unsure{\overset{s}{\rightarrow}|\overset{s}{\leftarrow}}$
$R^* \to C^*$ where $s$ is the unique multiplicative section gives a
canonical isomorphism $\mu_{n,k} \simeq \mu_{n,C} \simeq Z_n$ where the last isomorphism
comes from the distinguished nth root exp$(2\pi i/n)$. When $X = G_R$
then $\mu_{n,R} = \mu_{n,C} \simeq Z_n$ . Hence there is a canonical element
$C = \{C_n\} \in H^{2d}_X(E, \hat{Z})$ giving an isomorphism

(5) \[ \cup C : H^q(X, \hat{Z}) \xrightarrow{\sim} H^{q+2d}_X(E, \hat{Z}) . \]

Now choose a basepoint of $E - 0$ lying over the basepoint of $X$ .
I have been told that by careful analysis of the etale homotopy
type functor it is possible to rigidify $(E - 0)^\wedge_{et} \to E^\wedge_{et}$ to be-%
come a pro-object $\underline{f}$ in the homotopy category of maps in $C\underline{H}$ .
Let us assume that further analysis will show there is a canonical
isomorphism $H^*(C\underline{f}, \hat{Z}) \simeq H^*_X(E, \hat{Z})$ which is reasonable in virtue of
(4). Then (5) gives a Thom class $\unsure{U|U\cdot} \in H^{2d}(C\underline{f}, \hat{Z})$ showing that
$(\underline{f},u,U)$ is an oriented sphere fibration \$E with base $X^\wedge_{et}$ ,
where $u : \hat{X}_{et} \simeq \hat{E}_{et}$ is the isomorphism coming from the projection
$E \to X$ . As (5) is compatible with base-change $g^*(\$E) \simeq \$(g^*E)$
when $g$ is a map of preschemes.

\margin{\handwritten{64}}
